% ============================================================================
%  sections/A1_proofs.tex --- Appendix A: theory.
%
%  Organisation.  One subsection per object:
%    A.1  setup, notation, standing assumptions
%    A.2  the log-density estimator: change of variables, Hutchinson trace,
%         and the integration scheme (explicit Euler in the state, midpoint
%         quadrature for the divergence integral).  It targets log q at t=1
%         at EVERY step count; the step count controls accuracy, not target.
%    A.3  the flow-matching score read-out          (lem:score)
%    A.4  why the gradient condition is a trapped surrogate (prop:trap/conflict)
%    A.5  the value condition                       (thm:certificate, cor:gauge)
%    A.6  local-to-global identifiability           (thm:l2g, cor:basins)
%    A.7  the training objective as a surrogate for the metric (prop:surrogate)
%    A.8  scope of the theory
%
%  The force-gauge aside (lem:forcegauge / thm:gauge) was DELETED in the
%  2026-08-29 appendix-compression pass: no experiment exercised it, it was not
%  the explanation of the force result, and it was not a contribution.  Do not
%  re-add it.
%
%  Conventions that must not be reverted:
%    (1) free energy is A(c) = -kT log Z(c); F is ALWAYS the force.
%    (2) coordinates are x, never r.
%    (3) no ESS claim: we never measured ESS.
%    (4) the rigid-motion assumption is COM-removal only -- the code does NOT
%        quotient rotations.
%    (5) charge is threaded to the teacher, spin is not (singlet default).
%    (6) the L_value gradient is a frozen-trajectory / truncated-adjoint one,
%        so the theory is scoped to the zero set and identifiability of the
%        ideal residual condition and NOT to convergence of that update.
%    (7) the reported density is the CLAMPED POSITIONAL-FLOW density q_theta;
%        p_theta is reserved for the joint generator's conditional density and
%        appears only in ass:A2's contrast clause.
%
%  REPEATED CAVEATS ARE CONSOLIDATED.  Each of the following is stated once in
%  the main text and once here, at the anchor named, and is cross-referenced
%  rather than restated: q_theta is not the joint conditional (ass:A2); the
%  divergence is prior work (rem:logvar); the implemented gradient is a
%  truncated-adjoint surrogate (rem:idealresidual); the objective is local and
%  does not identify basin mass (cor:basins); kT = 1 eV is numerical
%  (app:hyper); the likelihood estimate is resolution-dependent (ass:A-disc,
%  app:estimator); value-arm means are conditional on completion (app:p0cells).
%
%  Theorem environments are declared in main.tex and the counters continue from
%  the main text, so appendix-only results start after Theorem 1 / Lemma 1 /
%  Corollary 1.  Proofs of main-text results use
%  \begin{proof}[Proof of \Cref{...}] and introduce no second numbering.
% ============================================================================

\section{Theory: Assumptions, Estimator, and Proofs}
\label{app:proofs}

This appendix collects everything the main text states without proof, in the order
the argument uses it: the standing assumptions (\Cref{app:setup}); what the
reported log-density actually is, derived from the flow (\Cref{app:density}); the
score read-out the gradient baseline relies on and the two reasons it is a poor
target (\Cref{app:score,app:trap}); the certificate behind the value objective and
exactly what it identifies (\Cref{app:certificate,app:l2g}); the identity linking
the training objective to the reported metric (\Cref{app:surrogate}); and a closing
inventory of which statements this paper exercises and which it only states
(\Cref{app:scope}).

\subsection{Setup, notation and standing assumptions}
\label{app:setup}

A molecular state is a pair $(c,\vx)$. Operationally the composition $c \in \mathcal{C}$ is the
element multiset with the atom count $N(c)$ it determines and the total charge
$q$, which is what the pipeline threads to the teacher, and $\mathcal{C}$ is
countable. Spin multiplicity $s$ is a standing restriction rather than a
component of $c$: it is never passed and the teacher defaults to a singlet
(\Cref{ass:charge}), so every statement below is at that default and the
composition audit used for identifiability keys on (element multiset, charge)
accordingly. The conformation
$\vx \in \R^{3N(c)}$ is the nuclear coordinate vector with the centre of mass
removed, so the effective dimension is $d = 3(N-1)$. The teacher is a universal
neural potential $E : (c,\vx) \mapsto \R$ \citep{omol25,esen},
smooth in $\vx$, with conservative forces $F(c,\vx) = -\nabla_{\vx} E(c,\vx)$,
and $kT>0$ is a fixed temperature scale in eV. Throughout, $F$ denotes the
\emph{force} and never a free energy.

The target is the joint Boltzmann law and its two factors,
\begin{equation}
  \pi(c,\vx) = \frac{e^{-E(c,\vx)/kT}}{\mathcal{Z}},
  \qquad
  Z(c) = \int e^{-E(c,\vx)/kT}\,\mathrm{d}\vx,
  \qquad
  A(c) = -kT \log Z(c),
  \label{eq:app-targets}
\end{equation}
so that $\mathcal{Z} = \sum_{c} Z(c) = \sum_c e^{-A(c)/kT}$,
$\pi(c) = e^{-A(c)/kT}/\mathcal{Z}$ and
$\pi(\vx \mid c) = e^{-E(c,\vx)/kT}/Z(c)$. Neither $Z(c)$ nor $\mathcal{Z}$ is
ever formed. The generative model is a flow $\mathrm{d}\vx_t/\mathrm{d}t =
v_\theta(\vx_t,t)$ transporting $p_0 = \mathcal{N}(0,\sigma^2 I)$ on the
centre-of-mass-free subspace at $t=0$ to $q_\theta := q_1^\theta$ at $t=1$ with
the discrete channels clamped at their endpoint labels (\Cref{ass:A2}),
trained on the linear interpolant
$\vx_t = (1-t)\vx_0 + t\vx_1$ \citep{lipman2023flowmatching,liu2023rectifiedflow,%
albergo2023interpolants}.

\begin{assumption}[Teacher regularity]
\label{ass:A1}
$E(c,\cdot)$ is the teacher potential at fixed $(N,q,s)$ and is $C^1$ in $\vx$ on
a connected domain.
\end{assumption}

\begin{assumption}[Which density is reported: the clamped positional-flow density]
\label{ass:A2}
The quantity every theorem below is stated for, and every experiment reports as
$\log q_\theta(\cdot \mid c)$, is the log-density of the \emph{clamped
positional flow}: the flow obtained by holding the discrete channels --- atom
types, formal charges and the empty bond tensor --- at their endpoint labels
$(a_1, c_1, e_1)$ for the whole reverse trajectory, while only $\vx$ evolves.
That object is positive on the sampled support and normalised on the
centre-of-mass-free subspace by construction, being an exact change of variables
applied to a normalised Gaussian prior \citep{ffjord}, up to the discretisation of
\Cref{ass:A-disc}. It coincides with the generator's own joint conditional density
$p_\theta(\cdot \mid c)$ only if the discrete trajectory is determined by its
endpoint; the generator evolves types and charges through a continuous-time Markov
chain, so in general the two differ, we have not quantified the difference, and
this is recorded as a limitation (\Cref{app:limitations}). The symbol $q_\theta$ is
used nowhere else in this appendix.
\end{assumption}

\begin{assumption}[Translation-invariant velocity field]
\label{ass:A3}
$v_\theta$ is invariant to global translation, so the ambient $3N$-dimensional
trace of its Jacobian equals the intrinsic $3(N-1)$-dimensional one
(\Cref{app:density}). The centre-of-mass-removed equivariant field of the backbone
satisfies this. \textbf{Rotations are not quotiented}: the prior log-density uses
effective dimension $3(N-1)$, i.e.\ centre-of-mass removal only. We do not claim a
rotation-quotiented density.
\end{assumption}

\begin{assumption}[Discretised likelihood]
\label{ass:A-disc}
The reported $\log q_\theta$ is a discretised continuous-flow likelihood ---
explicit Euler in the state, midpoint-rule quadrature for the divergence integral,
with a stochastic trace estimator --- so it is a biased and noisy estimate of the
exact clamped positional-flow log-density. The step count controls that bias; it
does not change which density is estimated. Step counts and probe counts are in
\Cref{app:hyper}; the consequences are worked through in
\Cref{app:density,app:estimator}.
\end{assumption}

\begin{remark}[What the theory is scoped to, given the implemented gradient]
\label{rem:idealresidual}
The gradient actually back-propagated through $\Ls_{\mathrm{value}}$ is a
truncated-adjoint surrogate: the reverse trajectory is advanced without gradient
tracking, each state is detached before its divergence is taken, and the prior term
is computed from a detached endpoint, so the parameter gradient flows only through
the divergence evaluations (\Cref{app:asimplemented}). We do not differentiate
through the ODE solve, and the implemented update is therefore not exact
likelihood-gradient descent. Every statement in this appendix is accordingly scoped
to the \emph{zero set and identifiability} of the ideal residual condition
$w \equiv \mathrm{const}$ --- what a model that had driven the population objective
to zero would have been shown, and what that leaves free --- and to nothing about
whether, or to what, the implemented update converges.
\end{remark}

\begin{assumption}[Rigid-motion and symmetry factors]
\label{ass:rigid}
$Z(c)$ in \eqref{eq:app-targets} is an integral over $\R^{3N}$ and therefore
carries translational and rotational volume factors and a symmetry-number
correction. We take $Z(c)$ to be the configurational integral in the
centre-of-mass-removed space that the density estimator actually models; as noted
in \Cref{ass:A3}, rotations are \emph{not} quotiented by the implementation.
Within a parent group this is benign: all $K$ perturbations are small
displacements of one data geometry in one frame, so any rotational volume
factor is approximately constant within the group and cancels from the per-group
correlation, and is exactly annihilated by the within-parent variance loss. It is
\emph{not} benign for any comparison across groups or across separated basins,
which is one more reason we make none (\Cref{thm:l2g}).
\end{assumption}

\begin{assumption}[Well-posed conditional energy]
\label{ass:charge}
$E(c,\cdot)$ is the potential at fixed $(N,q,s)$, so the teacher must be queried
with the correct charge and spin. Our implementation decodes the total charge from
the frozen charge channel and passes it to the teacher; the spin multiplicity is
\emph{not} passed and the teacher defaults to a singlet. For radicals and
open-shell complexes this is the wrong electronic state. Every reported evaluation
is on main-group organic parents drawn from the held-out split, where the singlet
default is the standard assignment for closed-shell species; we have not verified
the multiplicity of every evaluated parent. The restriction is recorded as a
limitation (\Cref{app:limitations}) and is the reason no radical, open-shell or
transition-metal result is reported.
\end{assumption}

\subsection{The reported log-density: change of variables, trace estimation, and the integration scheme}
\label{app:density}

Everything the value objective constrains, and everything the calibration
diagnostics measure, is the number this subsection defines. It is derived in three
steps --- an exact identity, an unbiased estimator of the one expensive term in it,
and an integration scheme --- and the third step is the one that decides how
\emph{accurately} the density is reported, not which density it is.

\paragraph{Step 1: the instantaneous change of variables.} Let $\vx_t$ solve
$\mathrm{d}\vx_t/\mathrm{d}t = v_\theta(\vx_t,t)$ and let $p_t$ be the law of
$\vx_t$. Mass conservation is the continuity equation
$\partial_t p_t + \nabla\!\cdot(p_t v_\theta) = 0$. Expanding the divergence,
$\nabla\!\cdot(p_t v_\theta) = \nabla p_t \cdot v_\theta + p_t \nabla\!\cdot v_\theta$,
and dividing by $p_t > 0$,
\begin{equation}
  \partial_t \log p_t + v_\theta \cdot \nabla \log p_t
  \;=\; -\,\nabla\!\cdot v_\theta .
\end{equation}
The left-hand side is the material derivative of $\log p_t$ along a trajectory, so
\begin{equation}
  \frac{\mathrm{d}}{\mathrm{d}t}\,\log p_t(\vx_t)
  \;=\; -\,\nabla\!\cdot v_\theta(\vx_t,t) ,
  \label{eq:app-icv}
\end{equation}
the instantaneous change of variables underlying continuous normalising flows
\citep{ffjord}. Integrating \eqref{eq:app-icv} from $t=0$ to $t=1$ along the
trajectory that ends at $\vx_1$ and using $p_0 = \mathcal{N}(0,\sigma^2 I)$ gives
\eqref{eq:ffjord} of the main text,
\begin{equation}
  \log q_\theta(\vx_1 \mid c)
  \;=\; \log p_0(\vx_0) - \int_0^1 \nabla\!\cdot v_\theta(\vx_t,t)\,\mathrm{d}t ,
  \qquad
  \vx_0 = \vx_1 - \int_0^1 v_\theta(\vx_t,t)\,\mathrm{d}t .
  \label{eq:app-ffjord}
\end{equation}
In practice the trajectory is generated \emph{backwards} from the data point
$\vx_1$, so one reverse solve returns both $\vx_0$ and the accumulated divergence.
\Cref{eq:app-ffjord} is exact in continuous time and requires no invertibility
argument beyond uniqueness of the ODE solution; every approximation in the reported
number enters in steps 2 and 3.

\paragraph{Step 2: the trace, and why a stochastic estimator is used.} The
divergence is the trace of the Jacobian, $\nabla\!\cdot v_\theta = \operatorname{tr} J$
with $J_{ij} = \partial (v_\theta)_i/\partial x_j$. Computing $\operatorname{tr} J$ exactly
costs $3N$ backward passes, one per coordinate, at every integration step. The
Hutchinson estimator \citep{hutchinson} replaces it by
\begin{equation}
  \operatorname{tr} J \;=\; \E_{\epsilon}\bigl[\epsilon^{\!\top} J \epsilon\bigr],
  \qquad \E[\epsilon] = 0,\quad \E[\epsilon\epsilon^{\!\top}] = I ,
  \label{eq:app-hutch}
\end{equation}
which is exact in expectation for any such $\epsilon$, and each sample costs one
vector--Jacobian product --- a single backward pass --- followed by an inner product
with $\epsilon$. We use $n_{\mathrm{p}} = 2$ Rademacher probes per step, i.e.\ two
backward passes in place of $3N$.

Two properties of \eqref{eq:app-hutch} matter for how the resulting numbers may be
read. First, with independent entries the estimator's variance is
\begin{equation}
  \Var\bigl[\epsilon^{\!\top} J \epsilon\bigr]
  \;=\; \sum_{i<j} \bigl(J_{ij}+J_{ji}\bigr)^{2}
  \;\;\xrightarrow{\;J = J^{\!\top}\;}\;\;
  2\Bigl(\lVert J\rVert_F^{2} - \textstyle\sum_i J_{ii}^{2}\Bigr) ,
  \label{eq:app-hutchvar}
\end{equation}
because $\epsilon^{\!\top} J\epsilon = \sum_i J_{ii} + \sum_{i<j}(J_{ij}+J_{ji})\epsilon_i\epsilon_j$
and the products $\epsilon_i\epsilon_j$ over distinct pairs are uncorrelated with unit
variance. Rademacher probes are the standard choice because the diagonal of $J$
contributes no variance at all, whereas Gaussian probes pay $2\lVert J\rVert_F^2$;
averaging $n_{\mathrm{p}}$ probes divides \eqref{eq:app-hutchvar} by
$n_{\mathrm{p}}$. Second, the divergence enters \eqref{eq:app-ffjord}
\emph{linearly}, so trace noise adds variance to the reported $\log q_\theta$ but
no bias: an unbiased trace at every step gives an unbiased estimate of the
discretised integral. The bias in the reported number comes from the quadrature,
not from the probes. We have not measured the repeat-to-repeat spread of
$\log q_\theta$ on the reported population, and say so as a limitation
(\Cref{app:limitations}).

Finally, \Cref{ass:A3} is what lets the probes be drawn in the ambient $3N$
coordinates without correcting the dimension. If $v_\theta$ is invariant to a
global translation then $J u = 0$ for each of the three uniform-translation
directions $u$, so those directions contribute zero eigenvalues and the ambient
trace equals the trace of $J$ restricted to the centre-of-mass-free subspace.
Rotations enjoy no such identity and are not quotiented anywhere in the
implementation.

\paragraph{Step 3: the integration scheme, and what the step count does and does not
control.} The reverse-time solve of \eqref{eq:app-ffjord} uses \emph{explicit Euler
in the state with midpoint-rule quadrature for the divergence integral}. With $n$
steps the implementation sets $\Delta t = 1/n$ and evaluates the field and its
divergence at
\begin{equation}
  t_k \;=\; 1 - \Bigl(k + \tfrac{1}{2}\Bigr)\Delta t ,
  \qquad k = 0,1,\dots,n-1 ,
  \label{eq:app-grid}
\end{equation}
accumulating $-\Delta t\,\nabla\!\cdot v_\theta(\vx_{t_k},t_k)$, advancing the state
by $\vx \leftarrow \vx - \Delta t\, v_\theta(\vx,t_k)$, and adding the prior
log-density at the end.

Two schemes are in play and they must not be conflated. The \emph{state} update is
explicit Euler with $n$ steps of size $1/n$, which traverses the whole of $[0,1]$
and carries $O(\Delta t)$ error. The \emph{divergence} integral is approximated by
the midpoint rule on $[0,1]$: all $n$ nodes are interior to the interval, the
interval itself is not truncated, and the rule carries $O(\Delta t^{2})$ error and
is exact for integrands linear in $t$. The half-step offset therefore avoids
evaluating the field at $t = 1$ without moving the endpoint the scheme targets. For
every $n$, what the estimator targets is $\log q_\theta$ at $t = 1$; there is no
family of smoothed marginals indexed by the step count, and refining $n$ estimates
the \emph{same} quantity more accurately rather than a different one. Analytic
checks on the implementation confirm this: on a field of constant divergence, and on
a time-varying one, the accumulated integral equals the full-interval value at every
$n$ rather than a truncated $1 - 1/(2n)$ value, and on a linear continuous
normalising flow the estimate converges to the closed-form endpoint log-density as
$n$ grows.

Step-count dependence in the reported numbers is accordingly \emph{numerical
discretisation error}, dominated by the $O(\Delta t)$ Euler state integration rather
than by the $O(\Delta t^{2})$ quadrature, together with whatever endpoint stiffness
the learned field carries. \Cref{app:estimator} tabulates the step counts, reports
what the absence of a plateau looks like on real molecules, and measures how much of
the reported effect size depends on the choice.

Two uses of the symbol $\epsilon$ must be kept apart, and neither is a smoothing
scale of the reported density: the $\epsilon$ of \eqref{eq:app-eps} below is the
Gaussian width $(1-t)\sigma/t$ attached to the flow-matching score read-out at a
probe time $t$ and not to the likelihood estimator, and it is unrelated to the
mass-relaxation $\varepsilon$ of \Cref{rem:epsilon}, which is written with a
different glyph throughout.

\subsection{The flow-matching score read-out}
\label{app:score}

The gradient-supervision baseline of \Cref{sec:method} needs the score of the
model's own time-$t$ marginal. Flow matching supplies it in closed form, which is
convenient and, as \Cref{app:trap} shows, also the source of the baseline's
trouble.

\begin{lemma}[Flow-matching score read-out]
\label{lem:score}
Let $x_t = (1-t)x_0 + t x_1$ with $x_0 \sim \mathcal{N}(0,\sigma^2 I)$ independent of
$x_1 \sim p_{\mathrm{data}}$ and $v^{*}(x,t) = \E[x_1 - x_0 \mid x_t = x]$. Then for every
$t \in [0,1)$, $\nabla_x \log p_t(x) = \bigl(t\,v^{*}(x,t) - x\bigr)/\bigl((1-t)\sigma^2\bigr)$.
\end{lemma}

\begin{proof}[Proof of \Cref{lem:score}]
Conditionally on the data endpoint, $\vx_t \mid \vx_1 \sim
\mathcal{N}(t\vx_1, (1-t)^2\sigma^2 I)$, so the time-$t$ marginal
$p_t(\vx) = \int \mathcal{N}(\vx; t\vx_1, (1-t)^2\sigma^2 I)\,p_{\mathrm{data}}(\vx_1)
\,\mathrm{d}\vx_1$ is a Gaussian mixture with covariance independent of the
mixing variable. Differentiating under the integral sign and using
$\nabla_{\vx}\mathcal{N}(\vx;\mu,\Sigma) = -\Sigma^{-1}(\vx-\mu)\mathcal{N}$ gives
the posterior-mean form
\begin{equation}
  \nabla_{\vx}\log p_t(\vx)
  = \frac{t\,\E[\vx_1 \mid \vx_t = \vx] - \vx}{(1-t)^2\sigma^2}.
  \label{eq:app-mixture-score}
\end{equation}
The flow-matching regression target is $\vx_1 - \vx_0$, and
$\vx_0 = (\vx_t - t\vx_1)/(1-t)$ gives $\vx_1 - \vx_0 = (\vx_1 - \vx_t)/(1-t)$.
Taking conditional expectations, the marginal velocity is
$v^\ast(\vx,t) = (\E[\vx_1\mid \vx_t=\vx] - \vx)/(1-t)$, i.e.
$\E[\vx_1 \mid \vx_t = \vx] = \vx + (1-t)v^\ast(\vx,t)$. Substituting into
\eqref{eq:app-mixture-score}, one factor of $(1-t)$ cancels and
\begin{equation}
  s^\ast(\vx_t,t) \;=\; \frac{t\,v^\ast(\vx_t,t) - \vx_t}{(1-t)\sigma^2}.
  \label{eq:app-score}
\end{equation}
Replacing $v^\ast$ by the learned $v_\theta$ defines $s_\theta$. Finally, the map
$v \mapsto s$ in \eqref{eq:app-score} is affine with linear part
$t/((1-t)\sigma^2)$, so a velocity perturbation $\delta v$ moves the implied
score by $t\|\delta v\|/((1-t)\sigma^2)$.
\end{proof}

\begin{remark}[Sign]
\label{rem:sign}
Matching $s_\theta \to +F/kT$ is Boltzmann-attracting because
$\nabla_{\vx}\log \pi = -\nabla_{\vx}E/kT = +F/kT$. The sign in
\eqref{eq:app-score} is not a convention that may be flipped.
\end{remark}

\begin{remark}[Amplification at the probe times used]
\label{rem:app-amplification}
With $\sigma = 1$, the factor $t/((1-t)\sigma^2)$ equals $5.7$, $11.5$ and $32.3$
at the three training probe times $t \in \{0.85, 0.92, 0.97\}$. The pole at
$t = 1$ is not a removable artefact: $s_\theta$ is the score of the
\emph{smoothed} marginal $p_t$, and the smoothing vanishes exactly as the
read-off amplification diverges.
\end{remark}

\subsection{Why the gradient condition is a trapped surrogate}
\label{app:trap}

Two mechanisms, both specific to matching \eqref{eq:app-score} \emph{off-policy at
data endpoints}, are compatible with the gradient baseline failing to help. Both
turn on one structural fact: the network is a function of the interpolation state
$\vx_t$ while the target is the force read at the data endpoint $\vx_1$. Both are
stated as \emph{diagnostic hypotheses} rather than as explanations of the
measurement in \Cref{sec:ablation}; neither is isolated experimentally, and neither
bears on force supervision in general \citep{psm,driftingboltzmann}.

\begin{proposition}[Bias--variance trap: a diagnostic hypothesis]
\label{prop:trap}
At the population optimum of $\Ls_{\mathrm{FM}}$, and with
$\epsilon = (1-t)\sigma/t$, the read-out is, up to a factor $1/t$, the score of
$p_{\mathrm{data}} * \mathcal{N}(0,\epsilon^2 I)$ evaluated at
$\vx_1 + \epsilon z$, $z \sim \mathcal{N}(0,I)$. Compared with the endpoint target
$F(\vx_1)/kT$ this incurs a scale error and a smoothing error, each $O(1-t)$ or
smaller, together with an evaluation-point displacement of size $\epsilon$, while
by \Cref{rem:app-amplification} the gradient noise injected into the parameters
grows as $(1-t)^{-2}$. No choice of $t$ makes the deterministic terms and the noise
term small together.
\end{proposition}

\begin{proposition}[Objective conflict: a diagnostic hypothesis]
\label{prop:conflict}
The population minimiser of $\Ls_{\mathrm{FM}}$ has implied score
$\nabla \log p^{\mathrm{data}}_t(\vx_t)$. The network input is $\vx_t$ while the
force target is the endpoint quantity $F(c,\vx_1)/kT = \nabla\log\pi(\vx_1)$, so the
population minimiser of the idealised \emph{squared} off-policy force loss is not
$\nabla\log\pi(\vx_1)$ but the conditional expectation
\begin{equation}
  s_\theta(\vx_t,t) \;\longrightarrow\;
  \E\bigl[\nabla\log\pi(\vx_1) \bigm| \vx_t, t, c\bigr]
  \;=\; \E\bigl[F(c,\vx_1)/kT \bigm| \vx_t, t, c\bigr],
  \label{eq:app-condexp}
\end{equation}
a smoothed and shrunk field. The two population targets differ whenever the corpus
is not drawn from the target law, $p_{\mathrm{data}} \neq \pi$, and the joint
minimiser is neither; a fixed $\lambda_1$ moreover weights the conflicting target
by $(1-t)^{-2}$.
\end{proposition}

\begin{proof}[Proof of \Cref{prop:trap} and \Cref{prop:conflict}]
Reparameterise the smoothing. Write $\vx_t = t\,y$ with
\begin{equation}
  y = \vx_1 + \epsilon z, \qquad z \sim \mathcal{N}(0,I), \qquad
  \epsilon = (1-t)\sigma/t ,
  \label{eq:app-eps}
\end{equation}
so that $p_t(\vx) = t^{-d} \tilde p_\epsilon(\vx/t)$ with
$\tilde p_\epsilon = p_{\mathrm{data}} \ast \mathcal{N}(0,\epsilon^2 I)$, and
therefore $s_\theta(\vx_t,t) = t^{-1}\nabla\log\tilde p_\epsilon(\vx_1+\epsilon z)$
at the population optimum. Against the target $\nabla\log\pi(\vx_1)$ this exposes
four mismatches: a $1/t$ prefactor contributing relative error $(1-t)/t = O(1-t)$;
a smoothing error $\nabla\log(p \ast \mathcal{N}_\epsilon) - \nabla\log p =
O(\epsilon^2) = O((1-t)^2)$ by the heat-flow expansion; an evaluation-point
displacement contributing $\epsilon\,\nabla^2\log\pi\,z$ to first order; and, since
by \Cref{lem:score} any velocity error is multiplied by $t/((1-t)\sigma^2)$ and the
loss is quadratic in the score, parameter-gradient noise scaling as $(1-t)^{-2}$.
The first three shrink as $t \to 1$ while the fourth diverges, so the total error is
bounded below by a trade-off of the form $O(1-t) + O((1-t)^{-2}\eta)$ with $\eta$
the velocity-estimation noise level, which is \Cref{prop:trap}; averaging over
probe times, as we do, reduces the constant but not the trade-off.

For \Cref{prop:conflict}, the displacement is not merely noise that averages away.
The loss is minimised over functions of $\vx_t$, and $(\vx_t,\vx_1)$ are dependent
but not mutually determined, so for a squared loss the population minimiser over
measurable functions of the input is the conditional expectation
\eqref{eq:app-condexp}, by the orthogonality property of $L^2$ projection; it equals
$\nabla\log\pi(\vx_1)$ only if $\vx_1$ is $\sigma(\vx_t)$-measurable, which the
interpolation makes false for every $t < 1$. The result is a shrinkage \emph{bias}
in the learned score rather than a variance term, and the two population targets ---
the score of the data interpolation marginal, and the conditional expectation of
endpoint Boltzmann scores --- are distinct whenever $p_{\mathrm{data}} \neq \pi$.

Two caveats bound both propositions. The implemented $\ell$ is Huberised and
per-atom norm-capped \eqref{eq:score}, so its population minimiser is a robustified
location functional of the same conditional law and not \eqref{eq:app-condexp}. And
both arguments concern population objectives at their exact optima, so they name
mechanisms the construction admits rather than ones we have isolated.
\end{proof}

\begin{remark}[The antecedent holds for our data, and on-policy variants are outside what we measure]
\label{rem:antecedent}
The training corpus is a curated union of relaxation trajectories, molecular
dynamics snapshots and conformer sets across many systems, charges and spins
\citep{omol25}. It is not an equilibrium sample from a single-temperature Boltzmann
distribution, and the training $kT = 1.0$~eV is not a physical temperature of the
dataset, so the antecedent $p_{\mathrm{data}} \neq \pi$ of \Cref{prop:conflict}
holds by construction. Querying forces at model samples rather than at data points
would partially repair the evaluation-point displacement and the objective conflict,
because the two targets would then refer to one distribution; it would not repair
the $(1-t)^{-1}$ amplification. No such variant is measured here: the
gradient-supervision cells of the matched grid (\Cref{app:p0cells}) carry the
off-policy endpoint term alone, whereas the earlier round's force cells also carried
an on-policy distillation term, which is why that round's \emph{force} comparisons
are confounded and are superseded by cell E (\Cref{app:cells}).
\end{remark}

\subsection{The value condition: a partition-function-free certificate}
\label{app:certificate}

The certificate below is the standard zero-of-the-log-variance-divergence
statement written in our conditional notation. Its structural property --- that an
additive constant in $\log\tilde\pi$, hence the partition function, cancels --- is
due to the authors cited in \Cref{rem:logvar} and is not claimed here. We reproduce
the proof because \Cref{cor:gauge} and \Cref{thm:l2g} both refer to its internals.

\begin{theorem}[Zero-variance certificate on a region]
\label{thm:certificate}
Under the regularity assumptions of \Cref{app:setup} (smooth $E$, positive normalised
$q_\theta$ on the COM-removed subspace, translation-invariant velocity field), fix a composition
$c$, a region $U$ and a reference measure $\rho$ with $\mathrm{supp}\,\rho = U$, and let
$w(x) = \log q_\theta(x \mid c) + E(c,x)/kT$. Then \emph{(i)} $\Var_{x\sim\rho}[w] = 0$ if and
only if $w$ is $\rho$-a.s.\ constant on $U$; if $\rho$ has a positive Lebesgue density on $U$ this
says $q_\theta(\cdot \mid c)$ restricted to $U$ is proportional to $e^{-E(c,\cdot)/kT}$, and if
$\rho$ is atomic on $K$ sampled points it says $\log q_\theta$ and $-E/kT$ differ by one common
constant at those $K$ points, i.e.\ $K-1$ scalar equations, with no statement about points outside
the sample; \emph{(ii)} if $\rho$ has a positive Lebesgue density on $U$ then the multiplier is
fixed \emph{relative to $U$}, $a = q_\theta(U \mid c)/Z_U(c)$ with
$Z_U(c) = \int_U e^{-E(c,x)/kT}\,\mathrm{d}x$ the integral over $U$ and not the global $Z(c)$, so
that $q_\theta(\cdot \mid c, U) = \pi(\cdot \mid c, U)$ exactly; if in addition
$q_\theta(U \mid c) = 1$ then $a = 1/Z_U(c)$, and $q_\theta(\cdot \mid c) = \pi(\cdot \mid c)$
globally only when $U$ also carries all of the target's mass, $Z_U(c) = Z(c)$;
\emph{(iii)} $\Var_{x\sim\rho}[w]$ is
invariant to $E \mapsto E + g(c)$ and needs no evaluation of $Z(c)$ or $\mathcal{Z}$.
\end{theorem}

\begin{corollary}[What the empirical loss identifies]
\label{cor:gauge}
The empirical objective \eqref{eq:lvalue} has null space exactly the functions that are
constant within each parent group. It therefore determines $\log q_\theta(\cdot \mid c)$ up to
one free additive constant \emph{per group}, and no further.
\end{corollary}

\begin{remark}[Attribution]
\label{rem:logvar}
$\Var_\rho[w]$ with $w = \log q_\theta - \log\tilde\pi$ and
$\tilde\pi = e^{-E/kT}$ is the log-variance divergence \citep{vargrad,
nusken2021pathspace,richter2024improvedsampling}. Part (iii) --- invariance to
an additive constant in $\log\tilde\pi$, hence no partition function --- is the
defining structural property of that divergence and is due to those authors, not
to this paper. Estimating it under a reference measure that is not the model's own
(off-policy) is likewise standard \citep{nusken2021pathspace,sendera2024offpolicy}.
\end{remark}

\begin{remark}[The certificate is about the exact clamped density; what is formed is its discretisation]
\label{rem:epscert}
\Cref{thm:certificate} is stated for the exact normalised $q_\theta$, whereas what
the objective and the metric both form is its Euler/midpoint estimate at a finite
step count (\Cref{app:density}). The residual the objective can enforce therefore
carries discretisation bias, which does not vanish at the step counts we can afford
(\Cref{app:estimator}); the target quantity is the same at every step count, and the
accuracy is not.
\end{remark}

\begin{proof}[Proof of \Cref{thm:certificate}]
Fix $c$, a region $U$ and $\rho$ with $\mathrm{supp}\,\rho = U$; let
$w(\vx) = \log q_\theta(\vx\mid c) + E(c,\vx)/kT$.

(i) A random variable has zero variance if and only if it is a.s.\ constant, so
$\Var_{\vx\sim\rho}[w] = 0$ iff $w \equiv \kappa$ $\rho$-a.s.\ on $U$, and then
$q_\theta(\vx\mid c) = a\,e^{-E(c,\vx)/kT}$ with the positive multiplier
$a = e^{\kappa}$, for $\rho$-almost every $\vx \in U$. Note the sign: the residual
constant enters the density through $e^{+\kappa}$, and the same holds component by
component in \Cref{thm:l2g}. The strength of that conclusion depends on $\rho$. If $\rho$ has a
positive Lebesgue density on $U$, the identity holds Lebesgue-almost everywhere on
$U$ and $q_\theta$ restricted to $U$ is proportional to the Gibbs weight. If
$\rho$ is the empirical measure on $K$ sampled points, the identity holds at those
$K$ points and nowhere else is constrained: the certificate is $K-1$ scalar
equations, not a density statement, and asserting proportionality on a region from
finitely many pointwise equations would be a strictly stronger claim than the
variance supplies. Conversely if $q_\theta(\cdot\mid c) \propto e^{-E(c,\cdot)/kT}$
$\rho$-a.e.\ on $U$ then $w$ is constant there. Note that (i) alone is a statement
\emph{about $U$}: nothing outside $U$ is constrained.

(ii) Let $\rho$ have a positive Lebesgue density on $U$, so that (i) gives the
identity Lebesgue-a.e.\ on $U$ with a positive multiplier $a = e^{\kappa}$. Since
$q_\theta(\cdot\mid c)$ is a normalised density (\Cref{ass:A2}), integrating
$a\,e^{-E/kT}$ over $U$ gives $q_\theta(U \mid c) = a\,Z_U(c)$, where
$Z_U(c) = \int_U e^{-E(c,\vx)/kT}\,\mathrm{d}\vx$ is the integral over $U$ and
\emph{not} the global $Z(c)$. Hence $a = q_\theta(U\mid c)/Z_U(c)$, and dividing,
the conditional law is pinned exactly:
$q_\theta(\vx \mid c, U) = e^{-E(c,\vx)/kT}/Z_U(c) = \pi(\vx \mid c, U)$. Suppose in
addition $q_\theta(U\mid c) = 1$, so that $a = 1/Z_U(c)$: this determines the
multiplier but identifies $q_\theta(\cdot\mid c)$ with $\pi(\cdot\mid c)$
\emph{globally} only when the target also puts all of its mass on $U$, that is when
$Z_U(c) = Z(c)$. Otherwise the model is the Gibbs law truncated to $U$ and
renormalised, differing from $\pi(\cdot\mid c)$ by the factor $Z(c)/Z_U(c)$ on $U$
and by all of $\pi(U^{\mathrm{c}}\mid c)$ outside it. It is exactly the proviso
$q_\theta(U\mid c) = 1$ that our per-parent perturbation groups violate, which is
what \Cref{thm:l2g} quantifies, and even where it holds the step to a global
statement needs the target's mass to sit on $U$ as well. For an atomic $\rho$ no
hypothesis of this form is available at all: a finite point set is Lebesgue-null,
so $q_\theta(U \mid c) = 0$ for any continuous model density and part (ii) is
inapplicable rather than merely unattainable. The $\varepsilon$-relaxation of
\Cref{rem:epsilon} repairs mass leaking outside a bounded region but not a
Lebesgue-null reference measure, and we make no part-(ii) claim in the empirical
reading at any $\varepsilon$.

(iii) $\Var_{\vx\sim\rho}[w]$ is unchanged by $E \mapsto E + g(c)$ for any
function $g$ of composition alone, because such a shift adds a constant to $w$
within the group; and it involves neither $Z(c)$ nor $\mathcal{Z}$, because the
within-group centring subtracts whatever constant those would contribute
(\Cref{rem:logvar}).
\end{proof}

\begin{proof}[Proof of \Cref{cor:gauge}]
The empirical objective is
$\Ls_{\mathrm{value}} = M^{-1}\sum_{m}\Var_k(w_{m,k})$ with
$w_{m,k} = \log q_\theta(\vx_{m,k}) + E_{m,k}/kT$, where $m$ indexes parent
groups and $k$ indexes the $K$ perturbations of parent $m$. A within-group
variance is annihilated exactly by functions that are constant within each group
and by nothing else, so the null space of $\Ls_{\mathrm{value}}$ is the set of
per-\emph{group} additive constants --- one scalar $b_m$ per group, not one per
composition. The objective therefore determines the \emph{shape} of
$\log q_\theta$ inside each group's support and no more. Fixing those constants
requires absolute energies, which is the role of an anchor term; the anchor was
inactive in every run reported in this paper (\Cref{app:asimplemented}). One
granularity caveat belongs with that statement: the anchor as implemented predicts
its constant from element counts, atom count and total charge, so it is keyed on
the \emph{composition} and can fix at most one constant per composition, not the
one constant per group this corollary identifies. Those coincide only when a
composition supplies a single group, which our training pool does not satisfy for
$31\%$ of its parents (\Cref{app:data}).
\end{proof}

\subsection{Local-to-global identifiability}
\label{app:l2g}

\Cref{thm:l2g} characterises the null space of the group-wise objective under
stated support and normalisation assumptions; connectivity of the overlap graph is
what removes the per-component offsets from that null space. It does not decide
whether a trained model is Boltzmann, and no reported run has a connected overlap
graph. Because the statement is sensitive to the choice of reference measure, we
name ours first and then separate the continuous and the empirical readings.


% ---------------------------------------------------------------------------
%  RELOCATED FROM THE MAIN TEXT BY THE 2026-08-19 COMPRESSION PASS.
%  Sec. 4 of the main text now carries a single short proposition
%  (prop:ident); the full statements live here, immediately before their
%  proofs.  Text is verbatim from the pre-compression main text.
% ---------------------------------------------------------------------------

\begin{definition}[Overlap graph]
\label{def:overlap}
Fix a composition $c$ with groups $j = 1,\dots,J$ and reference measures $\rho_j$ with supports
$U_j$, all dominated by a common $\sigma$-finite measure $\mu$, and write
$f_j = \mathrm{d}\rho_j/\mathrm{d}\mu$. The \emph{overlap graph} $\mathcal{G}_c$ has vertices
$\{1,\dots,J\}$ and an edge $\{j,j'\}$ if and only if
\begin{equation}
  \int \min\bigl(f_j, f_{j'}\bigr)\,\mathrm{d}\mu \;>\; 0 ,
  \label{eq:overlap}
\end{equation}
equivalently if and only if $\rho_j$ and $\rho_{j'}$ are \emph{not mutually singular}. The
criterion does not depend on the choice of $\mu$. An edge supplies the measure
$\mathrm{d}\nu = \min(f_j,f_{j'})\,\mathrm{d}\mu$, which is nonzero and absolutely continuous
with respect to both $\rho_j$ and $\rho_{j'}$; that is what forces the two constancy statements to
agree, and it is what a set merely charged by both measures does not supply. Overlap here is a
property of the measures and not of the topological supports: $U_j \cap U_{j'} \neq \emptyset$ is
necessary for an edge but far from sufficient.
\end{definition}

\begin{theorem}[Local-to-global identifiability]
\label{thm:l2g}
Let the population objective $\sum_j \pi_j \Var_{\rho_j}[w]$ vanish, so $w \equiv b_j$
$\rho_j$-a.s.\ on $U_j$ for each $j$, and let $U_c = \bigcup_j U_j$. Then \emph{(i)}
$b_j = b_{j'}$ for every edge of $\mathcal{G}_c$, so the constants agree within each connected
component. Suppose in addition that each $\rho_j$ has a positive Lebesgue density on $U_j$. Then
\emph{(ii)} if $\mathcal{G}_c$ is connected, $w$ is a single constant on $U_c$ and the model law
\emph{conditioned on} $U_c$ is the Gibbs law conditioned on $U_c$,
$q_\theta(x \mid c, U_c) = e^{-E(c,x)/kT}/Z_{U_c}(c)$ with
$Z_{U_c}(c) = \int_{U_c} e^{-E(c,x)/kT}\,\mathrm{d}x$; the global statement
$q_\theta(\cdot \mid c) = \pi(\cdot \mid c)$ follows only in the further case that $U_c$ carries
all of the mass of \emph{both} laws, $q_\theta(U_c \mid c) = 1$ and $Z_{U_c}(c) = Z(c)$;
\emph{(iii)} if $\mathcal{G}_c$ has $C \ge 2$ components with regions $U^{(1)},\dots,U^{(C)}$ and
$q_\theta(U_c \mid c) = 1$, the zero-loss densities supported on $U_c$ form the
$(C-1)$-parameter family carried by positive multipliers $a_i > 0$,
\begin{equation}
  q_\theta(x \mid c) = a_i\,e^{-E(c,x)/kT} \ \ (x \in U^{(i)}),
  \qquad \textstyle\sum_{i=1}^{C} a_i Z^{(i)} = 1, \quad
  Z^{(i)} = \int_{U^{(i)}} e^{-E(c,x)/kT}\,\mathrm{d}x ,
  \label{eq:l2g}
\end{equation}
in which every vector of occupancies
$\bigl(q_\theta(U^{(1)} \mid c),\dots,q_\theta(U^{(C)} \mid c)\bigr)
 = (a_1 Z^{(1)},\dots,a_C Z^{(C)})$
in the open simplex is attained: the \emph{relative} mass of separated basins is not identified.
\end{theorem}

Part (i) needs no density assumption and is proved below in two lines of measure
theory. Parts (ii) and (iii) are the continuous corollary, where a density and a
normalisation hypothesis are available; the empirical corollary
(\Cref{cor:empirical}) is the one our runs are in, and there parts (ii) and (iii)
have no counterpart, so what remains is an identifiability statement. Connectivity
is sufficient and, in that sense, necessary: a disconnected graph can still admit a
particular zero-loss solution whose constants happen to coincide, but only
connectivity forces a single constant for \emph{every} zero-loss solution.

\begin{corollary}[What our design identifies]
\label{cor:basins}
In every run reported here each parent contributes exactly one group, and $\rho_j$ is the empirical
measure on that group's sampled geometries: one data geometry together with its $K-1$
displacements. Distinct parents draw their geometries independently and continuously, so with
probability one no two groups share a sampled point, no edge is formed, and $\mathcal{G}_c$ is a
set of isolated vertices whatever $J$ is. Only part~(i) applies. The objective constrains the local
relation between $\log q_\theta$ and energy inside each cloud, is blind to that cloud's additive
offset, and constrains the density outside it, in particular the relative probability of
distinct basins, not at all.
\end{corollary}

The corollary rests on mutual singularity of the group measures and not on a
one-group-per-composition design, which our training pool does not have: $31\%$ of
its parents sit on a composition shared with another parent (\Cref{app:data}), and
those repeated clouds are still disconnected, which sharpens the point rather than
weakening it. Connectivity is therefore a property of the supervision design;
groups tiled along a path between basins would make $\mathcal{G}_c$ connected, which
is an experiment this paper does not contain. One further gap belongs here: the
objective and the evaluation metric average over \emph{different} clouds --- different
potentials, parents, displacement scales and estimator step counts
(\Cref{sec:setup}) --- so the corollary bounds what the objective can constrain and
does not identify the two.

\begin{remark}[Which reference measure is the operative one]
\label{rem:refmeasure}
What the objective \eqref{eq:lvalue} and the evaluation metric \eqref{eq:rm} both
average over is the \emph{empirical} measure on the sampled geometries of a
parent, so the empirical corollary below is the operative one throughout. Hence
$U_j$ is a finite point set, the certificate of \Cref{thm:certificate} is $K-1$
scalar equations rather than constancy on a region, and distinct parents have
almost surely mutually singular reference measures. Read instead as the full-support
Gaussian those points are drawn from, no two groups of a composition would be
mutually singular, every pair would therefore carry an edge, and part (ii) of
\Cref{thm:l2g} would apply --- but no run here approaches that population objective,
and reading our experiments through it would credit us with a certificate we do not
have. The empirical reading is the conservative one.
\end{remark}

\begin{corollary}[Empirical corollary: what zero loss certifies at sampled points]
\label{cor:empirical}
Let each $\rho_j$ be the empirical measure on the $K_j$ sampled geometries of group
$j$, and let the objective vanish. Then \emph{(a)} for each $j$, $w$ takes one
common value $b_j$ at the $K_j$ sampled points of group $j$, which is $K_j - 1$
scalar equations and constrains no unsampled point; \emph{(b)} taking $\mu$ to be
counting measure on the union of the sampled points, $f_j$ is group $j$'s atomic
weight function and $\int \min(f_j,f_{j'})\,\mathrm{d}\mu > 0$ holds exactly when
the two groups share at least one sampled geometry; for independently drawn
continuous displacements that has probability zero, so distinct groups are almost
surely mutually singular and $\mathcal{G}_c$ is edgeless; and \emph{(c)} consequently the objective places no constraint relating
the offsets $b_j$ of distinct groups, and the relative mass those groups receive is
unconstrained by the loss. Statement (c) is an identifiability statement and not a
normalisation decomposition: the family \eqref{eq:l2g} is derived only in the
continuous reading, where the $U^{(i)}$ carry positive Lebesgue measure.
\end{corollary}

\begin{proof}[Proof of \Cref{thm:l2g}]
By \Cref{thm:certificate}(i) applied to each group, vanishing of every summand
gives $w \equiv b_j$ $\rho_j$-a.s.\ on $U_j$.

(i) Let $\{j,j'\}$ be an edge of $\mathcal{G}_c$ (\Cref{def:overlap}), so with
$f_j = \mathrm{d}\rho_j/\mathrm{d}\mu$ the measure
$\mathrm{d}\nu = \min(f_j,f_{j'})\,\mathrm{d}\mu$ is nonzero. By construction
$\nu \le \rho_j$ and $\nu \le \rho_{j'}$ as measures, hence $\nu \ll \rho_j$ and
$\nu \ll \rho_{j'}$, so $w = b_j$ holds $\nu$-a.s.\ and $w = b_{j'}$ holds
$\nu$-a.s. The union of the two $\nu$-null exceptional sets is $\nu$-null and
$\nu \neq 0$, so there is a point at which both identities hold and
$b_j = b_{j'}$. This is precisely the step that the weaker ``some set charged by
both measures'' criterion does not license: two measures can each charge a common
set while being mutually singular, and then no single point need carry both
statements. Equality along edges propagates along paths, so $j \mapsto b_j$ is
constant on each connected component.

(ii) If $\mathcal{G}_c$ is connected, all $b_j$ equal a common $b(c)$, hence
$w \equiv b(c)$ on $U_c = \bigcup_j U_j$ and
$q_\theta(\vx\mid c) = a(c)\,e^{-E(c,\vx)/kT}$ there, with $a(c) = e^{b(c)} > 0$.
\Cref{thm:certificate}(ii) applies with $U = U_c$ and fixes the multiplier only
relative to $U_c$: $a(c) = q_\theta(U_c\mid c)/Z_{U_c}(c)$, so the conditional law
$q_\theta(\cdot \mid c, U_c)$ is exactly the Gibbs law conditioned on $U_c$. If in
addition $U_c$ carries all the model's mass, $q_\theta(U_c\mid c) = 1$, then
$a(c) = 1/Z_{U_c}(c)$; and only if $U_c$ also carries all the target's mass, so that
$Z_{U_c}(c) = Z(c)$, does this become $q_\theta(\cdot\mid c) = \pi(\cdot\mid c)$
globally. Neither hypothesis holds in any run reported here
(\Cref{rem:epsilon,cor:empirical}).

(iii) Work in the continuous reading, where by hypothesis each $\rho_j$ has a
positive Lebesgue density on $U_j$. Suppose $\mathcal{G}_c$ has components
$\mathcal{S}_1,\dots,\mathcal{S}_C$ with regions
$U^{(i)} = \bigcup_{j \in \mathcal{S}_i} U_j$. Two distinct
components have Lebesgue-null intersection: take $\mu = \lambda$ Lebesgue measure in
this reading, and suppose $U^{(i)} \cap U^{(i')}$ had positive Lebesgue measure.
Then some pair $j \in \mathcal{S}_i$, $j' \in \mathcal{S}_{i'}$ has
$\lambda(U_j \cap U_{j'}) > 0$, and on that common set $f_j$ and $f_{j'}$ are both
strictly positive by hypothesis, so $\min(f_j,f_{j'}) > 0$ there and
$\int \min(f_j,f_{j'})\,\mathrm{d}\lambda > 0$, which is \eqref{eq:overlap};
\Cref{def:overlap} would then place an edge between them, contradicting that they
lie in different components. Lebesgue-nullity is what the
normalisation decomposition below needs, which is why the positive-density
hypothesis is carried explicitly into parts (ii) and (iii) rather than built into
\Cref{def:overlap}: under the atomic reference measures our runs actually use,
every $U_j$ is Lebesgue-null and this decomposition is unavailable, which is why
\Cref{cor:empirical} states the weaker identifiability conclusion instead. By (i) the
residual is constant on each $U^{(i)}$, say $w \equiv b_i$, so
$q_\theta(\vx\mid c) = a_i\,e^{-E(c,\vx)/kT}$ on $U^{(i)}$ with $a_i = e^{b_i} > 0$;
conversely every such density lies in the zero-loss set of the objective, because
each $\Var_{\rho_j}[w]$ is then a variance of a constant. (This is a statement about
the objective's null set, not an existence claim about the model class: we do not
assert that a flow of the architecture used here realises every member of the
family.) Normalisation over $U_c$ is the single scalar constraint
$\sum_i a_i Z^{(i)} = 1$ on the positive multipliers, leaving $C-1$ free parameters,
which is \eqref{eq:l2g}. Finally $q_\theta(U^{(i)}\mid c) = a_i Z^{(i)}$, and as
$(a_1,\dots,a_C)$ ranges over the constraint set these occupancies range over the
whole open simplex; the Gibbs occupancies \emph{conditioned on $U_c$},
$Z^{(i)}/\sum_{i'}Z^{(i')}$, are one point of it and are singled out by nothing in
the objective. Note that these are conditional on $U_c$: the objective says nothing
about $\pi$'s mass outside $U_c$ either.
\end{proof}

\begin{proof}[Proof of \Cref{cor:basins}]
Each parent contributes exactly one group, and the operative reference measure is
the empirical measure on that group's sampled geometries (\Cref{rem:refmeasure}).
The conclusion follows from mutual singularity rather than from any count of groups
per composition. Two groups are joined in $\mathcal{G}_c$ exactly when their
empirical measures are not mutually singular, which for atomic measures means
exactly when the two clouds share a sampled geometry (\Cref{cor:empirical}(b));
distinct parents draw their displacements independently from continuous
distributions, so that event has probability zero and the two empirical measures
are almost surely mutually singular.
Every vertex is therefore isolated whatever $J$ is, each connected component is a
single group, and \Cref{cor:empirical} applies with one free offset per cloud.
\Cref{thm:certificate}(i) in its atomic form is the strongest part available.

Two remarks keep this honest. First, the argument is deliberately independent of
composition multiplicity, because the training pool does contain repeated
compositions (\Cref{app:data}); those repeated clouds are still edgeless, so the
conclusion is robust to multiplicity in a way that a $J = 1$ argument would not be.
Second, the conclusion concerns the supervision design and not the geometry of
Gaussian tails: were the $\rho_j$ read as full-support Gaussians, $\mathcal{G}_c$
would be complete, but then the population objective would be a variance over all of
$\R^{3N}$, which no run here approaches, and the \emph{measured} certificate would
still be the $K-1$ empirical equations. Even for a connected $\mathcal{G}_c$ the
normalisation hypothesis of \Cref{thm:l2g}(ii) is unattainable for a flow with
full-support prior, so the reachable continuous statement is the
$\varepsilon$-relaxed one of \Cref{rem:epsilon}.
\end{proof}

\begin{remark}[The normalisation hypothesis is an idealisation, and what survives without it]
\label{rem:epsilon}
No continuous normalising flow driven from a full-support Gaussian prior satisfies
$q_\theta(U_c \mid c) = 1$ for a bounded $U_c$, so parts (ii) of
\Cref{thm:certificate,thm:l2g} are limits rather than attainable states, and any upgrade
path through them must be read as approximate. The $\varepsilon$-relaxation is what
actually survives: if $w \equiv b$ on $U_c$ and $q_\theta(U_c \mid c) = 1 - \varepsilon$,
then with the positive multiplier $a = e^{b}$ we have $a\,Z_{U_c}(c) = 1-\varepsilon$, so
$a = (1-\varepsilon)/Z_{U_c}(c)$ and $b = -\log Z_{U_c}(c) + O(\varepsilon)$, and
the law of $q_\theta$ \emph{conditioned} on $U_c$ is exactly the Gibbs law conditioned on
$U_c$. Connectivity therefore buys a composition-level constant relative to the mass
inside $U_c$, not a global normalisation.
\end{remark}

\begin{remark}[What would test the theorem]
\label{rem:l2g-design}
\Cref{thm:l2g} is falsifiable: a model trained with $C > 1$ can have arbitrarily
wrong mode occupancies while achieving zero training loss and a per-group
correlation of $1$. Testing it requires groups that are not mutually singular ---
which, under the empirical reference measures of \Cref{rem:refmeasure}, means groups
that share sampled geometries: perturbation clouds tiled along a transition path so
that consecutive clouds re-use the same geometries, or groups drawn from one
reference ensemble obtained by enhanced sampling. Merely intersecting supports would
not suffice. That must be combined with a measurement of relative basin populations
rather than within-basin ordering. No experiment in this
paper does that (\Cref{sec:limitations}).
\end{remark}

\begin{remark}[Why cross-molecule pooling is broken]
\label{rem:crossbatch}
Across molecules the certificate holds with a \emph{molecule-dependent} constant
$-\log Z(c)$ whose spread is dominated by size and composition. A cross-batch
variance therefore measures $\log Z$ spread rather than Boltzmann deviation. This
is the same Simpson-type trap that forces the evaluation correlation to be computed
per parent molecule (\Cref{sec:setup}): one argument covers both the loss and the
metric, and it is why the within-parent form is the only one used in either.
\end{remark}

\subsection{Algebraic relation between residual variance and ordering}
\label{app:surrogate}

The objective penalises a within-group variance while the endpoint reports a
within-group correlation. They share one residual form, which is why optimising the
first is expected to move the second --- under a stated condition that our own
measurements show is not met.

\begin{proposition}[Residual-variance / correlation identity]
\label{prop:surrogate}
Fix a parent with $K$ evaluation geometries; let $u_k = -E_k/kT$ be the independent
evaluation energy, $L_k = \log q_\theta(x_k)$, $w_k = L_k - u_k$,
$\tau = \mathrm{sd}_k(w)/\mathrm{sd}_k(u)$ and $\rho = \mathrm{corr}_k(w,u)$. Then the
per-group Pearson correlation is $r = (1+\rho\tau)/\sqrt{1+2\rho\tau+\tau^2}$, reducing at
$\rho = 0$ to $r = 1/\sqrt{1+\tau^2}$, and that parent's contribution to \eqref{eq:lvalue}
is exactly $\tau^2\Var_k(u)$.
\end{proposition}

\begin{proof}[Proof of \Cref{prop:surrogate}]
Fix a parent $m$ with $K$ evaluation perturbations. Write $u_k = -E_k/kT$ for the
independent evaluation energy, $L_k = \log q_\theta(\vx_k)$, and
$w_k = L_k - u_k$ for the Boltzmann residual, and set
$\tau = \mathrm{sd}_k(w)/\mathrm{sd}_k(u)$ and $\rho = \mathrm{corr}_k(w,u)$.
From $L = u + w$,
\begin{equation}
  \Cov(L,u) = \Var(u)\,(1+\rho\tau),
  \qquad
  \Var(L) = \Var(u)\,(1 + 2\rho\tau + \tau^2),
\end{equation}
so the per-group Pearson correlation is
\begin{equation}
  r_m = \frac{1+\rho\tau}{\sqrt{1+2\rho\tau+\tau^2}},
  \qquad\text{and, if } \rho = 0, \qquad
  r_m = \frac{1}{\sqrt{1+\tau^2}} .
  \label{eq:app-surrogate}
\end{equation}
At $\rho = 0$ the right-hand side is strictly decreasing in $\tau$, and
$\Ls_{\mathrm{value}}$ for parent $m$ is exactly $\tau^2\Var_k(u)$. Hence at
$\rho = 0$ minimising the training loss minimises the residual that the reported
metric penalises. Away from $\rho = 0$ the two can move in opposite directions,
which is what \Cref{sec:scale} measures: the observed disagreement shows that the
simplifying condition does not hold, without identifying the sign or the
distribution of $\rho$ (\Cref{rem:nrvtau}). This proposition is a purely algebraic
identity, not a guarantee that reducing the loss improves the correlation.
\end{proof}

\begin{remark}[$\tau^2$ \emph{is} the normalised residual variance]
\label{rem:nrvtau}
With $u_k = -E_k/kT$ we have $\Var_k(u) = \Var_k(E/kT)$, so the metric
\eqref{eq:rm} satisfies
\begin{equation}
  \mathrm{NRV}_m \;=\; \frac{\Var_k[w]}{\Var_k[E/kT]} \;=\; \tau_m^{2},
  \label{eq:app-nrvtau}
\end{equation}
and by \Cref{prop:surrogate} parent $m$'s contribution to \eqref{eq:lvalue} is
$\mathrm{NRV}_m\,\Var_k(u)$. Within a single group under a single potential, then,
$\mathrm{NRV}$ is exactly the $\tau^2$ of \Cref{prop:surrogate}, and the metric is
the normalised evaluation analogue of the same Boltzmann-residual form the training
objective uses. The two are not one quantity in different units: they share that
residual structure but are evaluated with different potentials (eSEN against
GFN2-xTB), different populations (training shards against held-out parents) and at
different discretisation accuracies ($n = 4$ against $n = 12$; the same target
quantity, different Euler bias). The identity is also the
bridge between the ordering result and
the calibration result, and it shows that the two cannot both be read through the
$\rho = 0$ idealisation. At $\rho = 0$, $r = 1/\sqrt{1+\tau^2}$ is strictly
decreasing in $\tau$, so the arm with the larger $r$ must have the smaller
$\mathrm{NRV}$. The measurement is the opposite: value supervision has both the
larger $r$ ($0.369$ against $0.200$) and the larger $\mathrm{NRV}$ ($2.144$ against
$1.402$). These two endpoints are aggregated differently --- $r$ is a mean over
parents and $\mathrm{NRV}$ a median over parents (\Cref{app:scalediag}) --- so no
per-parent $\rho$ can be recovered by combining them. What the disagreement does
establish is that the zero-correlation simplification is not an adequate explanation
of the observed arm-level behaviour: under it the arm with the larger $r$ must have
the smaller $\mathrm{NRV}$, and no aggregation convention makes both reported
endpoints consistent with that. The sign and the distribution of the residual
correlation are \emph{not identified} by these aggregates, and we report neither.
\end{remark}

\begin{remark}[The correlation certifies shape, not temperature or calibration]
\label{rem:affine}
Pearson $r$ is invariant under positive affine maps of either variable, so it
cannot distinguish $p \propto e^{-E/kT}$ from $p \propto e^{-E/kT'}$ for any
$T' > 0$, and it is equally blind to the \emph{magnitude} of the log-density's
response. The per-group regression slope, the scale ratio
$\mathrm{sd}_k[\log q_\theta]/\mathrm{sd}_k[E/kT]$, the implied effective
temperature $kT/\mathrm{slope}$ and the normalised residual variance
\eqref{eq:rm} are the read-outs that are \emph{not} affine-invariant, which is why
they are reported alongside every correlation (\Cref{app:scalediag}). Their verdict
differs from the correlation's, and that difference is the paper's main empirical
statement.
\end{remark}

\begin{remark}[Attenuation, and why we do not lean on it]
\label{rem:attenuation}
By \Cref{ass:A-disc} the evaluated $\log q_\theta$ carries discretisation bias and
stochastic trace noise. The classical attenuation model,
$r_{\mathrm{meas}} \approx r_{\mathrm{true}}/\sqrt{1 + \Var(\text{noise})/\Var(\log p)}$,
would make the reported correlation an under-estimate of the true one, and
geometry-dependent discretisation bias would enter $\Var_k(w)$ as an estimator
floor. This is an argument and not a measurement, and we do not lean on it --- the
more so because what we \emph{can} measure points the other way: the step counts we
can afford show no plateau, and the gap shrinks precisely as the estimator becomes
more accurate, which is what a partly artefactual gap at coarse resolution would
also do (\Cref{app:estimator}). What is measured instead is the teacher-evaluator agreement
reference: pushing the teacher's own energies through the identical pipeline gives
$r = 0.927$ on the primary population, so a model reproducing its teacher exactly
would score about that rather than $1$, and our $0.369$ is about $40\%$ of it in
absolute terms (\Cref{app:scalediag}). Whatever accounts for the residual, it is
not teacher--evaluator disagreement.
\end{remark}

\subsection{Scope of the theory}
\label{app:scope}

Four inventories, so that no statement above is credited with more than it does.

\paragraph{Not ours.} \Cref{thm:certificate} and its partition-function-free
property are the log-variance divergence of
\citet{vargrad,nusken2021pathspace,richter2024improvedsampling}, restated in our
conditional notation; off-policy estimation of such a divergence is likewise
established \citep{sendera2024offpolicy}. See \Cref{rem:logvar}.

\paragraph{Exercised empirically.} \Cref{thm:certificate}(i) in its atomic form,
applied to each disjoint sampled cloud --- that is the training objective; and
\Cref{rem:crossbatch}, which dictates the per-group construction of both the loss
and the metric.

\paragraph{Stated as a diagnostic identity, with its hypothesis contradicted.}
\Cref{prop:surrogate} is exact algebra, but its $\rho = 0$ specialisation --- under
which reducing normalised residual variance must improve the correlation --- is
contradicted by the measurement of \Cref{sec:scale}, where the value arm has both
the larger $r$ and the larger $\mathrm{NRV}$ (\Cref{rem:nrvtau}). That failure is
itself one of the paper's findings.

\paragraph{Stated and \emph{not} tested here.} \Cref{thm:certificate}(ii) and
\Cref{thm:l2g}(ii): no reported run has a connected overlap graph, so we never
reach the regime in which the certificate becomes a statement about $q_\theta$ on a
region --- and even there it would pin only the law \emph{conditioned} on that
region (\Cref{rem:epsilon}). \Cref{thm:l2g}(iii) is the regime we \emph{are} in,
and its failure mode --- arbitrary relative basin weights at zero loss --- is
untested in both directions, because no cross-basin measurement exists in this
paper (\Cref{rem:l2g-design}, \Cref{sec:limitations}).
\Cref{prop:trap,prop:conflict} are candidate mechanisms for the negative result of
\Cref{sec:ablation}, consistent with it but not isolated experimentally.

\paragraph{Outside the theory altogether.} Nothing above predicts the
\emph{magnitude} of the correlations we report, and nothing above speaks to the
\emph{calibration} of $\log q_\theta$: \Cref{thm:certificate}(i) fixes the residual
only up to a per-group constant, and by \Cref{rem:affine} the per-group correlation
is invariant to positive affine maps of either variable. The scale-sensitive
read-outs --- normalised residual variance, regression slope, scale ratio and
effective temperature --- are therefore empirical additions with no counterpart
here, and the finding that they fail to improve while the correlation does
(\Cref{sec:scale}) is not something the theory anticipates; the only bridge between
the two is the bookkeeping identity \eqref{eq:app-nrvtau}. Finally, every statement
above concerns population objectives at their exact optimum, whereas
\eqref{eq:lvalue} is estimated from $M$ parents with a discretised, stochastic
$\log q_\theta$, and the update we run is the truncated-adjoint surrogate of
\Cref{rem:idealresidual}. Everything above is therefore a statement about the zero
set of the ideal residual condition and about what that zero set leaves
unidentified; nothing above is a statement about whether, or to what, the
implemented update converges.
